\chapter{Sacroiliac Joint Injection}

\section*{Overview}
A sacroiliac joint injection places local anesthetic, often with a corticosteroid, into or around the sacroiliac joint to help diagnose and treat pain arising from that joint.

\section*{Why It Is Done}
It may be considered for persistent lower-back, buttock, or leg pain when examination suggests sacroiliac joint involvement and conservative treatment has not provided adequate relief.

\section*{How It Is Performed}
The patient lies face down. After the skin is cleaned and numbed, the clinician guides a needle to the joint using appropriate procedural guidance and injects the medication. The patient is observed briefly afterward and may record pain relief over the following hours and days.

\section*{Preparation}
Tell the clinician about blood thinners, diabetes, infection, allergies, pregnancy, and previous reactions to injection medicines. Medicines may need adjustment only under medical direction. Arrange help with transportation if sedation is used.

\section*{Risks and Recovery}
Temporary soreness, bleeding, infection, nerve irritation, and a short-lived increase in pain are possible. Corticosteroids may temporarily affect blood glucose and other steroid-sensitive conditions. Most people resume light activity shortly afterward, following the clinician's instructions.

\section*{Outcome and Follow-up}
Short-term meaningful pain relief supports the possibility that the sacroiliac joint contributes to symptoms. Lack of relief suggests that another pain source may need evaluation. Repeated injections are limited according to clinical response and safety considerations.

\section*{References}
\begin{enumerate}
  \item MedlinePlus. Sacroiliac joint injection. U.S. National Library of Medicine; 2025.
  \item Current Medical Diagnosis and Treatment. McGraw-Hill; 2025.
\end{enumerate}
\clearpage
